% =====================================================================
% Bab II — Tinjauan Pustaka  [REVISI]
% Perubahan signifikan dan referensi baru ditandai \textcolor{blue}{...}
% =====================================================================
\chapter{Tinjauan Pustaka}
\label{bab:tinjauan-pustaka}

Bab ini membahas tinjauan terhadap literatur yang mendasari penelitian,
\textcolor{blue}{disusun dalam delapan sub bagian:
(1)~perawatan prediktif dalam konteks Industri 4.0,
(2)~vibrasi \emph{rolling element bearing},
(3)~diagnostik \emph{rolling element bearing} berbasis \emph{deep learning},
(4)~\emph{explainable AI} (XAI) pada tingkat atribusi masukan,
(5)~prognostik \emph{rolling element bearing} dan estimasi sisa umur pakai,
(6)~interpretabilitas mekanistik berbasis \emph{sparse autoencoder},
(7)~XAI pada sinyal satu dimensi, dan
(8)~Kesenjangan literatur yang memposisikan penelitian ini relatif
terhadap penelitian sebelumnya.}
Hasil kesenjangan yang teridentifikasi
pada §\ref{sec:bab2_gap_literatur} menjadi \textcolor{blue}{justifikasi langsung bagi
ketiga rumusan masalah} yang dirumuskan pada \textcolor{blue}{§\ref{sec:bab1_rumusan}}.

% Catatan perubahan struktur:
% - §II.6 (fisika vibrasi) DIPINDAHKAN menjadi §II.2 agar frekuensi
%   karakteristik sudah terdefinisi sebelum digunakan di §II.3 (diagnostik).
% - §II.7 baru ditambahkan: tinjauan XAI lintas domain.
% - Urutan seksi lama digeser satu angka ke bawah.

% -----------------------------------------------------------------
\section{Perawatan Prediktif dan Industri 4.0}
\label{sec:bab2_pdm_industri}
% -----------------------------------------------------------------

Strategi perawatan mesin industri dapat ita bagi menjadi tiga periode, dimana hal ini  
mencerminkan perkembangan kemampuan pengukuran, komputasi, dan konektivitas.
Pemahaman terhadap landasan konseptual ini penting untuk mengapresiasi
posisi perawatan prediktif sebagai kerangka yang menyatukan akuisisi data kondisi
secara kontinu dengan model prediksi umur mesin.

\subsection{Evolusi Perawatan: Korektif, Preventif, dan Prediktif}
\label{subsec:bab2_evolusi_perawatan}

Perawatan korektif (reaktif) menunggu terjadinya kegagalan sebelum dilakukan
tindakan perbaikan. Pendekatan ini berisiko tinggi karena kegagalan yang tidak
terduga dapat menghentikan keseluruhan proses produksi dan mengakibatkan kerugian
yang besar. Sedangkan perawatan preventif menjadwalkan penggantian komponen berdasarkan
interval waktu yang tetap, namun sering menghasilkan penggantian prematur saat
komponen masih layak digunakan, sehingga menambah biaya tanpa manfaat yang nyata
\citetitb{Mobley2002PlantMaintenance}.

Perawatan prediktif (\emph{predictive maintenance}, PdM), yang
distandardisasikan dalam ISO~13374 untuk akuisisi dan pemrosesan data kondisi
serta ISO~13381 untuk estimasi sisa umur, mengeliminasi kedua kelemahan
tersebut dengan mengintervensi berdasarkan kondisi aktual komponen
\citetitb{Carvalho2019PdMReview}. Dalam paradigma PdM, \emph{Condition-Based
Maintenance} (CBM) mendefinisikan tindakan perawatan yang dipicu secara langsung
oleh informasi kondisi yang diperoleh dari sensor. \emph{Prognostics and
Health Management} (PHM) memperluas CBM dengan menambahkan kemampuan prediksi:
bukan hanya mendeteksi degradasi yang sedang berlangsung, tetapi juga
memperkirakan kapan bagian mesin akan mencapai batas kegagalan (\emph{remaining
useful life}, RUL). Kemampuan prediksi ini memungkinkan perencanaan
perawatan yang lebih efisien, pengurangan \emph{unplanned downtime}, dan
optimasi persediaan suku cadang \citetitb{Lei2018PrognosticsReview}.

\emph{Rolling element bearing} merupakan objek utama penerapan PHM karena perannya yang
kritis pada emua mesin rotasi. Lebih dari 40--50\% kegagalan mesin
industri berakar pada kerusakan \emph{bearing} \citetitb{Lei2018Machinery}.
Mekanisme kegagalan \emph{bearing} mencakup \emph{spalling} pada permukaan lintasan
dalam (\emph{inner race}) atau luar (\emph{outer race}), keausan elemen
bergulir (\emph{ball} atau \emph{roller}), dan retakan sangkar (\emph{cage}).
Setiap mekanisme menghasilkan tanda vibrasi yang berbeda, yang menjadi dasar
untuk diagnostik berbasis sinyal.

\subsection{Ekosistem IoT dan Konteks \emph{Making Indonesia} 4.0}
\label{subsec:bab2_iot_industri4}

Inisiatif \emph{Making Indonesia} 4.0 yang diluncurkan pada April 2018
menetapkan transformasi digital pada lima sektor manufaktur prioritas yaitu makanan
dan minuman, tekstil, kimia, otomotif, dan elektronika
\citetitb{Kemenperin2018Making4}. Penerapan teknologi \emph{Internet of Things}
(IoT), analitik \emph{big data}, dan kecerdasan artifisial ditetapkan sebagai
pendorong utama peningkatan produktivitas dan daya saing
\citetitb{WEF2018Industry4}.

Dalam konteks ini, ekosistem sensor industri berkembang pesat menuju
arsitektur \emph{edge--cloud} berlapis. Pada lapis pertama (\emph{edge IoT}),
sensor akselerometer, suhu, dan kecepatan dipasang langsung pada mesin dan
dihubungkan ke \emph{gateway} industri seperti SKF~IMx-8 yang mampu melakukan
pra-pemrosesan sinyal lokal sebelum mengirimkan data ke lapisan di atasnya
\citetitb{SKF2017TechReport}. Pada lapis kedua (\emph{edge server}), komputasi
yang lebih intensif, seperti klasifikasi kondisi dan deteksi anomali, dilakukan di
\emph{server} lokal pabrik. Pada lapis ketiga (awan/\emph{cloud}), estimasi
RUL dan analisis historis dilakukan dengan sumber daya GPU yang lebih besar
\citetitb{Cheng2019EdgeCloud}. Arsitektur tiga-tier ini menjadi kerangka
konseptual yang menyatukan komponen diagnostik dan prognostik dalam penelitian
ini, sebagaimana diuraikan pada §VI.2.

Penerapan PdM berbasis IoT pada sistem produksi PT~SKF Indonesia (mesin
\emph{grinding} OR1/OR2 dengan akuisisi melalui SKF~IMx-8 dan penyimpanan di
AWS Cloud) menjadi konteks industri yang menajdi motivasi dalam penelitian disertasi ini. Integrasi
penuh data kualitas heterogen (vibrasi dan \emph{radial clearance checking})
sebagai masukan multi-modal merupakan \emph{roadmap} jangka panjang yang
ditetapkan sebagai \emph{future work} pada §VI.5.

% -----------------------------------------------------------------
% §II.2 BARU — dipindahkan dari §II.6 lama
% -----------------------------------------------------------------
\textcolor{blue}{\section{Vibrasi \emph{Rolling Element Bearing}}
\label{sec:bab2_fisika_vibrasi}}

\textcolor{blue}{Pemahaman terhadap fisika vibrasi \emph{rolling element bearing} memberikan \emph{ground
truth} yang diperlukan untuk (i)~memvalidasi apakah pola atribusi SHAP DeepExplainer
pada Bab~IV sesuai dengan impuls transien fisik, dan (ii)~memvalidasi pemetaan
SAE $\to$ BPFx pada Bab~V.
Bagian ini merangkum derivasi frekuensi karakteristik \emph{bearing} dan prosedur
analisis \emph{envelope} yang digunakan untuk mengukur amplitudo spektrum pada
tiap frekuensi target.
Dengan meletakkan dasar fisika sebelum bagian diagnostik (§\ref{sec:bab2_diagnostik})
dan prognostik (§\ref{sec:bab2_prognostik}), pembaca dapat mengikuti terminologi
BPFx yang digunakan secara konsisten pada bab-bab selanjutnya.}

\subsection{Frekuensi Karakteristik \emph{Bearing}}
\label{subsec:bab2_bpfx}

Mekanisme kegagalan \emph{bearing} menghasilkan rangkaian impuls periodik dengan
frekuensi yang bergantung pada geometri \emph{bearing} dan kecepatan poros. Untuk
\emph{bearing} dengan $n$ elemen bergulir, diameter elemen $d_b$, diameter
\emph{pitch circle} $d_m$, sudut kontak $\alpha$, dan kecepatan poros $f_r$
(dalam Hz), empat frekuensi karakteristik utama (\emph{Ball Pass Frequencies}
dan turunannya) diturunkan dari geometri kinematik sebagai berikut
\citetitb{McFadden1984BearingModel}:

\begin{equation}
  f_{\text{BPFO}} = \frac{n}{2}\, f_r \left(1 - \frac{d_b}{d_m}\cos\alpha\right)
  \label{eq:bab2_bpfo}
\end{equation}

\begin{equation}
  f_{\text{BPFI}} = \frac{n}{2}\, f_r \left(1 + \frac{d_b}{d_m}\cos\alpha\right)
  \label{eq:bab2_bpfi}
\end{equation}

\begin{equation}
  f_{\text{BSF}} = \frac{d_m}{2\,d_b}\, f_r
    \left[1 - \left(\frac{d_b}{d_m}\cos\alpha\right)^{\!2}\right]
  \label{eq:bab2_bsf}
\end{equation}

\begin{equation}
  f_{\text{FTF}} = \frac{f_r}{2} \left(1 - \frac{d_b}{d_m}\cos\alpha\right)
  \label{eq:bab2_ftf}
\end{equation}

\emph{Ball Pass Frequency Outer} ($f_\text{BPFO}$, Persamaan~\eqref{eq:bab2_bpfo})
merupakan frekuensi impuls yang ditimbulkan saat elemen bergulir melewati
cacat pada lintasan luar (\emph{outer race}). Karena lintasan luar diam
(tidak berputar) pada sebagian besar aplikasi, BPFO langsung berkaitan dengan
frekuensi elemen bergulir melewati titik cacat. \emph{Ball Pass Frequency
Inner} ($f_\text{BPFI}$, Persamaan~\eqref{eq:bab2_bpfi}) adalah analog untuk
cacat pada lintasan dalam; nilai BPFI selalu lebih tinggi dari BPFO karena
lintasan dalam berputar bersama poros. \emph{Ball Spin Frequency}
($f_\text{BSF}$, Persamaan~\eqref{eq:bab2_bsf}) merupakan frekuensi rotasi
elemen bergulir terhadap porosnya sendiri; cacat pada elemen bergulir
menghasilkan impuls pada frekuensi ini. \emph{Fundamental Train Frequency}
($f_\text{FTF}$, Persamaan~\eqref{eq:bab2_ftf}) merupakan frekuensi rotasi
sangkar (\emph{cage}); cacat pada sangkar terekspresikan pada harmonik FTF
\citetitb{Randall2011CondMonitoring}.

\textcolor{blue}{Untuk \emph{bearing} CWRU ($n = 8$ bola, $\alpha = 15{,}17^{\circ}$,
$d_b = 0{,}3126$~inci, $d_m = 1{,}537$~inci) pada kecepatan putar 1.750~rpm,
nilai teoritis adalah: $f_\text{BPFO} \approx 104{,}6$~Hz,
$f_\text{BPFI} \approx 157{,}9$~Hz, $f_\text{BSF} \approx 68{,}9$~Hz,
$f_\text{FTF} \approx 13{,}1$~Hz \citetitb{Smith2015CWRU}.}
Frekuensi-frekuensi ini
dan harmoniknya menjadi \emph{ground truth} fisika yang digunakan untuk
memvalidasi pemetaan SAE pada §V.3 dan untuk menginterpretasikan pola impuls
transien pada FSM di §IV.6.
\textcolor{blue}{Nilai-nilai ini konsisten dengan geometri dan kecepatan
yang dilaporkan pada §IV.11.}

\subsection{Analisis \emph{Envelope} dan Kurtogram Spektral}
\label{subsec:bab2_envelope}

Sinyal vibrasi \emph{bearing} yang rusak $x(t)$ merupakan superposisi komponen
periodik deterministik, modulasi amplitudo akibat kegagalan, dan komponen
acak. Analisis \emph{envelope} mengisolasi amplitudo modulasi untuk
mengungkapkan frekuensi karakteristik yang tersembunyi dalam komponen
\emph{carrier} frekuensi tinggi \citetitb{McFadden1984BearingModel}.

Prosedurnya: (1)~\emph{band-pass filtering} pada pita frekuensi resonansi
dominan, (2)~transformasi Hilbert untuk mendapatkan sinyal analitik
$z(t) = x_\text{bp}(t) + j\,\hat{x}_\text{bp}(t)$ dengan
$\hat{x}_\text{bp}(t)$ sebagai transformasi Hilbert dari sinyal ter-\emph{filter},
(3)~ekstraksi selubung $A(t) = |z(t)|$, dan (4)~komputasi spektrum selubung
$A_\varphi = \int_{[\varphi-\Delta,\, \varphi+\Delta]} |\mathcal{F}\{A(t)\}(f)|\,df$
untuk $\varphi \in \{\text{BPFO}, \text{BPFI}, \text{BSF}, \text{FTF}\}$.

Kurtogram spektral (\emph{spectral kurtosis}) \citetitb{Antoni2007Kurtogram}
digunakan untuk pemilihan pita \emph{band-pass} secara otomatis: pita yang
memaksimalkan kurtosis spektral mengandung komponen impulsif yang paling kuat,
dan karenanya paling informatif untuk analisis \emph{envelope}. Dalam konteks
pemetaan SAE, amplitudo $A_{r,\varphi}$ pada rekaman ke-$r$ dikorelasikan
dengan aktivasi fitur SAE $z_{r,i}$ menggunakan korelasi Pearson untuk mengukur
derajat korespondensi antara representasi laten dan fisika degradasi.

% -----------------------------------------------------------------
\section{Diagnostik \emph{Rolling Element Bearing} Berbasis \emph{Deep Learning}}
\label{sec:bab2_diagnostik}
% -----------------------------------------------------------------

Diagnostik \emph{bearing} bertujuan mengklasifikasikan kondisi komponen ke dalam
kategori kerusakan yang telah ditentukan, berdasarkan fitur yang diekstrak
dari sinyal vibrasi. Perkembangan pendekatan diagnostik berlangsung dari
algoritma berbasis model dengan ekstraksi fitur manual menuju arsitektur yang belajar
representasi secara langsung dari sinyal mentah (\emph{raw signal}).

\subsection{Evolusi Arsitektur Diagnostik}
\label{subsec:bab2_evolusi_arsitektur}

Pendekatan konvensional mengandalkan fitur domain waktu dan frekuensi yang
dihitung secara manual, yaitu nilai efektif (RMS), kurtosis, faktor puncak,
dan \emph{spectral centroid}, yang kemudian diumpankan ke pengklasifikasi
seperti \emph{Support Vector Machine} (SVM). Kelemahan utama pendekatan ini
adalah ketergantungan pada keahlian domain dalam rekayasa fitur dan
kesulitan dalam menangkap interaksi fitur yang kompleks \citetitb{Lei2018Machinery}.

Kemajuan \emph{convolutional neural network} (CNN) membuka era baru, sinyal
mentah menjadi masukan langsung, dan model belajar representasi secara
otomatis melalu lapisan yang mendalam. \citenameitb{Guo2016ADCNN} menunjukkan bahwa CNN satu dimensi mampu
mengklasifikasikan kondisi \emph{bearing} CWRU dengan akurasi tinggi tanpa rekayasa
fitur manual. \citenameitb{He2016ResNet} memperkenalkan koneksi residual yang
memungkinkan jaringan sangat dalam dilatih tanpa masalah \emph{vanishing gradient},
dan \citenameitb{Wen2020ResNet50} mengadaptasi ResNet-50 untuk diagnostik
\emph{bearing} dengan preprocessing \emph{continuous wavelet transform} menjadi gambar
dua dimensi.

Arsitektur \emph{Wide Deep CNN} (WDCNN) yang diperkenalkan oleh
\citenameitb{Zhang2017WDCNN} menjadi tonggak penting dalam diagnostik berbasis
sinyal mentah. Kernel pertama yang berukuran lebar (64 sampel, \emph{stride} 16)
dirancang secara eksplisit untuk ekstraksi fitur yang peka frekuensi, meniru
karakteristik filter \emph{bandpass} yang selektif terhadap komponen frekuensi
karakteristik \emph{bearing}. Kernel-kernel berikutnya (ukuran 3) menangkap pola
lokal yang lebih halus, menghasilkan kombinasi ekstraksi fitur \emph{global}
dan \emph{lokal} yang efektif.

Perkembangan lebih lanjut mencakup arsitektur berbasis Transformer dan
mekanisme \emph{attention} untuk diagnostik, yang menawarkan fleksibilitas
tangkap konteks global. Namun kecenderungan terkini menunjukkan bahwa akurasi
diagnostik pada dataset benchmark sudah mendekati saturasi di atas 99\%; fokus
penelitian bergeser ke aspek \emph{trustworthiness}, yakni kemampuan model untuk
menjelaskan keputusannya kepada teknisi perawatan
\citetitb{Serradilla2022DLPdM}.

\textcolor{blue}{\subsection{\emph{Transfer Learning} dan \emph{Domain Adaptation}: Tren dan Posisi Penelitian}
\label{subsec:bab2_transfer_learning}}

\textcolor{blue}{Sejak 2020, \emph{transfer learning} dan \emph{domain adaptation} telah menjadi
arah dominan dalam diagnostik \emph{bearing} berbasis \emph{deep learning}.
Pendekatan ini melatih model pada \emph{source domain} (kondisi operasi atau
jenis \emph{bearing} tertentu) lalu mengadaptasinya ke \emph{target domain}
yang berbeda dengan data berlabel terbatas \citetitb{Li2020TransferReview}.
Metode seperti \emph{domain adversarial neural networks} (DANN)
\citetitb{Ganin2015DANN} dan \emph{maximum mean discrepancy} (MMD) secara
aktif mengurangi perbedaan distribusi antar domain dalam ruang representasi.}

\textcolor{blue}{Meskipun demikian, penelitian ini tidak mengambil arah \emph{transfer learning}
dengan pertimbangan berikut: (i)~pertanyaan penelitian berfokus pada
\emph{interpretabilitas} keputusan model, bukan pada generalisasi lintas domain;
(ii)~protokol \emph{benchmark} multi-model pada dataset CWRU yang sama membutuhkan
kondisi evaluasi yang terkontrol dan seragam; dan (iii)~integrasi metode XAI
dengan model yang telah di-\emph{fine-tune} melalui domain adaptation
menambahkan lapisan kompleksitas yang berada di luar cakupan disertasi ini.
Perluasan kerangka FSM ke skenario \emph{cross-domain} ditetapkan sebagai
\emph{future work} pada §VI.5.}

\subsection{Dataset CWRU dan Isu \emph{Temporal Leakage}}
\label{subsec:bab2_cwru}

Dataset \emph{Case Western Reserve University} (CWRU) merupakan tolak ukur (\emph{benchmark})
paling banyak digunakan dalam penelitian diagnostik \emph{rolling element bearing}
\citetitb{Smith2015CWRU}. \emph{Test rig} terdiri dari motor 2~HP, kopling
torsi, dan \emph{bearing} SKF~6205-2RS~JEM yang dipasang pada sisi \emph{drive-end}
(DE) dan sisi \emph{fan-end} (FE). Cacat (\emph{defect}) artifisial disisipkan menggunakan
\emph{electro-discharge machining} (EDM) pada tiga komponen dengan tiga tingkat
keparahan (0,007; 0,014; 0,021 inci diameter cacat), menghasilkan 10~kelas
kondisi: Normal, tiga kelas \emph{inner race} (IR), tiga kelas \emph{outer
race} (OR), dan tiga kelas \emph{ball fault} (Ball) \citetitb{Smith2015CWRU}.

Sinyal DE diakuisisi pada 48~kHz di bawah empat kondisi beban motor (Load~0
hingga Load~3). Segmentasi dilakukan dengan \emph{window} 2.048 titik, menghasilkan
sampel yang cukup panjang untuk menangkap setidaknya satu siklus degradasi penuh
pada frekuensi karakteristik terendah (FTF $\approx$ 11,9~Hz pada 1.797~rpm →
periode 84~ms → 4.032 sampel; satu \emph{window} cukup untuk dua impuls FTF).

Kasus \emph{temporal leakage} sering terlewat dari perhatian dalam literatur awal,
pembagian acak (\emph{random split}) antara data latih dan uji menyebabkan
sampel yang berdekatan secara temporal, yang sangat berkorelasi satu sama lain,
terdistribusi di kedua himpunan. Hal ini menghasilkan estimasi akurasi yang
berlebihan karena model ``mengingat'' pola yang telah dilihat, bukan benar-benar
mengeneralisasi ke data baru \citetitb{Neupane2020CWRUReview}. Protokol
\emph{temporal split} (54/13/33\% = 1.240/310/750 sampel, berturutan dalam
waktu) digunakan dalam penelitian ini untuk menghindari kebocoran ini.

\subsection{Klasifikasi \emph{Bearing} dengan Multi-Model}
\label{subsec:bab2_multi_model}

Penelitian ini mengevaluasi tiga kelompok algoritma klasifikasi secara paralel
pada dataset CWRU. Ketiganya tidak dikombinasikan menjadi satu prediktor tunggal,
melainkan dijalankan masing-masing sebagai \emph{benchmark} untuk memahami
\emph{trade-off} antara akurasi, kompleksitas, keterdeplian di \emph{edge},
dan kedalaman interpretabilitas.

\textbf{Model berbasis kernel (SVM dan LR).} \emph{Support Vector Machine} dengan
kernel RBF dan \emph{Logistic Regression} \emph{one-vs-rest} beroperasi pada
18 fitur rekayasa manual: 9 fitur domain waktu (RMS, kurtosis, faktor puncak,
\emph{peak-to-peak}, kemiringan, standar deviasi, rata-rata, variansi, MAD) dan
9 fitur domain frekuensi (sentroid spektral, entropi spektral, energi pita
BPFO/BPFI/BSF/FTF, dan turunannya). SHAP \emph{KernelExplainer} memungkinkan
atribusi fitur yang bersifat model-agnostik untuk kedua model ini
\citetitb{TotoSuharto2024Conf1SVM}.

\textbf{Model berbasis pohon keputusan (\emph{tree-based model}).} \emph{Decision Tree} (DT),
\emph{Random Forest} (RF), dan XGBoost mengoperasikan pada himpunan fitur yang
sama. XGBoost dengan \emph{gradient boosting} dan regularisasi L1/L2
menghasilkan akurasi tertinggi dalam kelompok ini; RF memberikan \emph{ensemble}
bagging yang lebih stabil dibanding DT tunggal. SHAP \emph{TreeExplainer}
memberikan nilai Shapley yang eksak dan efisien untuk seluruh model berbasis pohon
\citetitb{TotoSuharto2024Conf2Tree}.

\textbf{Model berbasis Deep Learning(\emph{deep learning}).} WDCNN beroperasi langsung pada
sinyal mentah 2.048 sampel, tanpa proses rekayasa fitur secara manual. SHAP \emph{DeepExplainer}
menghasilkan nilai atribusi per posisi waktu (2.048 nilai per sampel per kelas),
yang kemudian diagregasi menjadi \emph{Fault Signature Maps} (FSM) sebagai
\emph{fingerprint} kelas kerusakan \citetitb{TotoSuharto2025Journal1FSM}.

\textcolor{blue}{Sebagai perbandingan dengan literatur sebelumnya,
Tabel~\ref{tab:bab2_klasifikasi_survey} menampilkan representasi penelitian dari tiga model berbeda.}
Tabel~\ref{tab:bab2_klasifikasi_survey} merangkum kinerja dan kemampuan
interpretabilitas dari representasi tiap model pendekatan. \textcolor{blue}{Angka kunci
dari penelitian ini} dilaporkan dari
\citetitb{TotoSuharto2024Conf1SVM}, \citetitb{TotoSuharto2024Conf2Tree},
dan \citetitb{TotoSuharto2025Journal1FSM}; detail lengkap dijelaskan pada
§IV.7--§IV.9.

\begin{table}[ht]
  \centering
  \caption{Ringkasan kinerja dan kemampuan interpretabilitas tiga model pendekatan
           algoritma klasifikasi \emph{bearing} pada dataset CWRU\textcolor{blue}{,
           dengan baris perbandingan dari literatur eksternal ditandai~($*$).}}
  \label{tab:bab2_klasifikasi_survey}
  \footnotesize
  \begin{tabularx}{\textwidth}{@{}l l l c l c@{}}
    \toprule
    \textbf{Algoritma} & \textbf{Kelompok} & \textbf{Dataset}
      & \textbf{Akurasi (\%)} & \textbf{Varian SHAP} & \textbf{FSM} \\
    \midrule
    SVM RBF      & Berbasis kernel  & CWRU & 96,4  & \emph{KernelExplainer} & — \\
    LR (\emph{one-vs-rest}) & Berbasis kernel & CWRU & 94,3 & \emph{KernelExplainer} & — \\
    \textcolor{blue}{SVM-RBF (Santos dkk., 2021)$^*$} & \textcolor{blue}{Berbasis kernel} & \textcolor{blue}{CWRU} & \textcolor{blue}{95,1} & \textcolor{blue}{\emph{KernelExplainer}} & \textcolor{blue}{—} \\
    RF            & Berbasis pohon  & CWRU & 95,8  & \emph{TreeExplainer}  & — \\
    XGBoost       & Berbasis pohon  & CWRU & 97,1  & \emph{TreeExplainer}  & — \\
    \textcolor{blue}{XGBoost (Cerrada dkk., 2022)$^*$} & \textcolor{blue}{Berbasis pohon} & \textcolor{blue}{CWRU} & \textcolor{blue}{96,8} & \textcolor{blue}{\emph{TreeExplainer}} & \textcolor{blue}{—} \\
    WDCNN (var A) & Deep Learning & CWRU & 99,87 & \emph{DeepExplainer}  & \checkmark \\
    WDCNN (var C, tanpa BN) & Deep Learning& CWRU & \textcolor{blue}{96,13} & \emph{DeepExplainer} & \checkmark \\
    \textcolor{blue}{CNN-1D \citetitb{Zhang2017WDCNN}$^*$} & \textcolor{blue}{Jaringan dalam} & \textcolor{blue}{CWRU} & \textcolor{blue}{99,60} & \textcolor{blue}{—} & \textcolor{blue}{—} \\
    \bottomrule
  \end{tabularx}

  \smallskip
  {\footnotesize Sumber: \citenameitb{TotoSuharto2024Conf1SVM} (SVM/LR);
   \citenameitb{TotoSuharto2024Conf2Tree} (RF, XGBoost);
   \citenameitb{TotoSuharto2025Journal1FSM} (WDCNN). Var~A = \emph{baseline}
   (kernel 64 + BN); Var~C = tanpa \emph{BatchNorm}. Detail hiperparameter
   di Lampiran E dan F.
   \textcolor{blue}{($^*$) Baris perbandingan dari literatur eksternal;
   protokol evaluasi masing-masing studi dapat berbeda dengan protokol
   \emph{temporal split} yang digunakan dalam penelitian ini.}}
\end{table}

Tabel~\ref{tab:bab2_klasifikasi_survey} menunjukkan perkembangan yang jelas.
Model kernel dan model berbasis pohon memberikan akurasi yang layak untuk \emph{edge
deployment} dengan interpretabilitas di tingkat fitur, sedangkan WDCNN
melampaui semua \emph{baseline} dalam akurasi sekaligus memungkinkan
interpretabilitas di tingkat sinyal mentah melalui FSM. Temuan bahwa
\textcolor{blue}{penonaktifan \emph{BatchNorm} (Varian~C) meningkatkan
\emph{discriminability} FSM dari 0,315 menjadi 0,735 pada subset ablasi
(30 \emph{background} / 50 sampel uji) --- kenaikan 0,420 poin
atau sekitar +133\% ($2{,}3\times$) relatif terhadap Varian~A dengan BN aktif} --- merupakan
kontribusi pertama yang dilaporkan dalam literatur diagnostik \emph{bearing}
(\emph{novelti N2}).

% -----------------------------------------------------------------
\section{\emph{Explainable AI}}
\label{sec:bab2_xai}
% -----------------------------------------------------------------

\emph{Explainable AI} (XAI) bertujuan menjelaskan alasan di balik keputusan
model sehingga dapat diverifikasi oleh operator manusia. Dalam konteks diagnostik
\emph{bearing}, XAI yang baik harus mampu menjelaskan \emph{mengapa} model menyimpulkan
suatu jenis kerusakan, serta idealnya mengkonfirmasi atau mengoreksi intuisi
rekayasawan yang berpengalaman dalam \emph{vibration analysis}.

\subsection{Taksonomi Metode XAI}
\label{subsec:bab2_taksonomi_xai}

Metode XAI untuk jaringan saraf dalam dapat dikelompokkan ke dalam tiga kategori
berdasarkan mekanisme kerjanya. Kategori pertama adalah \textbf{metode berbasis
gradien}: \emph{Saliency maps} \citetitb{Simonyan2013Saliency} menggunakan
gradien model terhadap masukan sebagai ukuran kepentingan; Grad-CAM
\citetitb{Selvaraju2017GradCAM} melokalisasi aktivasi kelas pada peta fitur
konvolusi. Kelemahan utama metode ini adalah sensitivitasnya terhadap
\emph{noise} masukan dan ketidakstabilan antar prediksi yang serupa.
\emph{Integrated Gradients} (IG) \citetitb{Sundararajan2017IG} mengatasi
masalah ini dengan mengintegrasikan gradien sepanjang jalur dari \emph{baseline}
ke masukan aktual, memenuhi aksiom \emph{sensitivity} dan \emph{implementation
invariance}.

Kategori kedua adalah \textbf{metode berbasis propagasi}: DeepLIFT
\citetitb{Shrikumar2017DeepLIFT} mempropagasikan kontribusi mundur berdasarkan
perbedaan aktivasi terhadap referensi (\emph{rescale rule}); LRP
(\emph{Layer-wise Relevance Propagation}) mengalokasikan relevansi dari lapisan
output ke masukan mengikuti aturan konservasi.

Kategori ketiga adalah \textbf{metode berbasis perturbasi}: LIME
\citetitb{Ribeiro2016LIME} memfit model linier lokal pada sekitar prediksi
yang dijelaskan; SHAP \citetitb{Lundberg2017SHAP} menghitung nilai Shapley
dari teori permainan kooperatif yang memenuhi serangkaian aksiom keadilan
yang ketat. Untuk domain diagnostik \emph{bearing}, metode berbasis perturbasi, khususnya SHAP,
paling luas diadopsi karena jaminan aksiomatiknya dan dukungan implementasi
untuk berbagai jenis model \citetitb{Bengherbia2024SHAPBearing}.

\subsection{SHAP \emph{DeepExplainer}}
\label{subsec:bab2_shap_deep}

SHAP \citetitb{Lundberg2017SHAP} memberikan kerangka interpretabilitas yang
berlandaskan teori permainan kooperatif: nilai Shapley $\phi_i$ untuk fitur
ke-$i$ mengukur kontribusi marginal rata-rata fitur tersebut terhadap prediksi
model di atas semua kemungkinan koalisi fitur lainnya. Sifat-sifat aksiomatik
yang dipenuhi, yaitu \emph{local accuracy}, \emph{missingness}, dan
\emph{consistency}, menjamin bahwa atribusi SHAP koheren dan tidak
menyesatkan secara matematis.

\emph{DeepExplainer} mengimplementasikan SHAP untuk jaringan saraf dalam
menggunakan pendekatan berbasis DeepLIFT (\emph{rescale rule}) dari
\citenameitb{Shrikumar2017DeepLIFT}: atribusi dihitung sebagai perbedaan
tertimbang aktivasi model terhadap aktivasi pada sampel \emph{background}
referensi, dipropagasikan mundur melalui lapisan-lapisan jaringan. Untuk
WDCNN dengan masukan sinyal mentah 2.048 sampel, \emph{DeepExplainer}
menghasilkan nilai atribusi $\phi^{(c)}(x)[p]$ untuk setiap posisi waktu
$p \in \{1, \ldots, 2048\}$ pada setiap sampel $x$ terhadap kelas target $c$.
Hasilnya adalah $500 \times 2048 \times 10$ nilai atribusi (500 sampel uji,
2.048 posisi, 10 kelas), yang menjadi bahan baku pembentukan FSM
\citetitb{TotoSuharto2025Journal1FSM}.

\subsection{SHAP per Jenis Model: Kapan Menggunakan yang Mana}
\label{subsec:bab2_shap_per_jenis}

Pemilihan varian SHAP yang tepat bergantung pada jenis model yang
dijelaskan. Ketiga varian memiliki karakteristik komputasi, keakuratan, dan
cakupan model yang berbeda secara fundamental, sehingga menggunakan varian
yang tidak sesuai dapat menghasilkan atribusi yang tidak akurat atau tidak
efisien.

\textbf{\emph{KernelExplainer}} bersifat model-agnostik: bekerja untuk model
apa pun yang dapat dievaluasi pada masukan yang diberikan. Pendekatan ini
memperkirakan nilai Shapley melalui perturbasi berbasis \emph{kernel}. Biayanya
tinggi, dengan kompleksitas $\mathcal{O}(TK^2)$ terhadap ukuran \emph{background}
$K$ dan jumlah fitur $T$, sehingga hanya praktis digunakan dengan
\emph{background} kecil. Dalam penelitian ini, \emph{KernelExplainer}
digunakan untuk SVM dan LR karena kedua model tidak memiliki struktur internal
yang dapat dieksploitasi oleh varian yang lebih spesifik
\citetitb{TotoSuharto2024Conf1SVM}.

\textbf{\emph{TreeExplainer}} mengimplementasikan algoritma TreeSHAP yang
menghitung nilai Shapley secara eksak dan efisien untuk model berbasis pohon
melalui traversal pohon polinomial \citetitb{Lundberg2017SHAP}. Tidak
memerlukan \emph{background} dataset; hasilnya adalah nilai Shapley yang
tepat (bukan aproksimasi). Dalam penelitian ini, \emph{TreeExplainer}
digunakan untuk DT, RF, dan XGBoost \citetitb{TotoSuharto2024Conf2Tree}.

\textbf{\emph{DeepExplainer}} mengeksploitasi struktur jaringan saraf dalam
melalui aturan \emph{rescale} DeepLIFT. Memerlukan \emph{background} 300 sampel
terstratifikasi (30 per kelas dari set latih); target eksplanasi 500 sampel uji
(50 per kelas). Kompleksitas bergantung pada ukuran \emph{background} dan
dimensi masukan; untuk masukan 2.048 titik dengan \emph{background} 300, komputasi
berlangsung dalam hitungan menit pada GPU \citetitb{TotoSuharto2025Journal1FSM}.

Tabel~\ref{tab:bab2_shap_variants} merangkum karakteristik ketiga varian.

\begin{table}[ht]
  \centering
  \caption{Perbandingan tiga varian SHAP \emph{explainer} yang digunakan dalam
           penelitian ini. \emph{Background} = dataset referensi yang diperlukan
           untuk menghitung nilai dasar.}
  \label{tab:bab2_shap_variants}
  \footnotesize
  \begin{tabularx}{\textwidth}{@{}l l l X l@{}}
    \toprule
    \textbf{Varian} & \textbf{Tipe} & \textbf{Akurasi} & \textbf{Target model}
      & \textbf{\emph{Background}} \\
    \midrule
    \emph{KernelExplainer}
      & Model-agnostik & Aproksimasi & Model apa pun (SVM, LR, dll.)
      & Diperlukan (sampling acak) \\
    \emph{TreeExplainer}
      & Eksak (TreeSHAP) & Eksak & DT, RF, XGBoost
      & Tidak diperlukan \\
    \emph{DeepExplainer}
      & DeepLIFT-\emph{based} & Aproksimasi & Jaringan saraf dalam
      & Diperlukan ($n = 300$) \\
    \bottomrule
  \end{tabularx}
\end{table}

\subsection{\emph{Fault Signature Maps} (FSM)}
\label{subsec:bab2_fsm}

\emph{Fault Signature Maps} (FSM) merupakan kontribusi dari
\citetitb{TotoSuharto2025Journal1FSM}: agregasi posisional nilai SHAP
\emph{DeepExplainer} yang menghasilkan \emph{fingerprint} per kelas kerusakan
dalam domain waktu resolusi penuh (2.048 posisi). Sementara SHAP \emph{pointwise}
menghasilkan atribusi spesifik per sampel (berguna untuk menjelaskan satu
prediksi), FSM meringkas perilaku model secara \emph{global} untuk seluruh
kelas kerusakan, menghasilkan peta yang stabil dan dapat dibandingkan
antar kelas.

Tiga varian FSM didefinisikan berdasarkan cara agregasi nilai SHAP:
\textbf{FSM Bertanda} (\emph{Signed}) merata-ratakan nilai atribusi dengan
arah aslinya, menunjukkan apakah posisi tertentu cenderung memberikan bukti
positif atau negatif untuk kelas yang bersangkutan; \textbf{FSM Mutlak}
(\emph{Absolute}) merata-ratakan nilai absolut, menghasilkan peta kepentingan
tanpa informasi arah; varian ini merupakan pilihan utama untuk membentuk
\emph{fingerprint}; \textbf{FSM Variansi} (\emph{Variance}) menghitung variansi nilai atribusi
antar sampel untuk mengukur stabilitas posisi sebagai indikator kelas.

Validitas FSM diukur melalui tiga metrik formal (definisi kuantitatif
disajikan pada §IV.6): (i)~\emph{discriminability index}: sejauh mana
FSM antar-kelas dapat dibedakan; (ii)~\emph{severity monotonicity fraction}:
fraksi posisi yang kepentingannya meningkat monoton seiring tingkat
keparahan cacat; (iii)~\emph{split-half stability coefficient}: \textcolor{blue}{kemiripan
kosinus (\emph{cosine similarity})} antara FSM yang dihitung dari dua subset acak separuh data uji.
\textcolor{blue}{Pada WDCNN Varian~A, koefisien stabilitas mencapai 0,940
pada analisis utama (300 \emph{background} / 500 sampel uji),
mengkonfirmasi bahwa FSM bukan artefak \emph{sampling}, melainkan representasi
yang stabil dari mekanisme keputusan model. Pada subset ablasi tereduksi
(30 \emph{background} / 50 sampel uji), stabilitas Varian~A tercatat 0,705
dan sedikit meningkat menjadi 0,741 pada Varian~C tanpa \emph{BatchNorm}
\citetitb{TotoSuharto2025Journal1FSM}.}

% -----------------------------------------------------------------
\section{Prognostik \emph{Rolling Element Bearing}: Estimasi Sisa Umur Pakai}
\label{sec:bab2_prognostik}
% -----------------------------------------------------------------

Prognostik \emph{bearing} bertujuan memperkirakan waktu yang tersisa sebelum
komponen mencapai kondisi kegagalan (\emph{remaining useful life}, RUL).
Berbeda dengan diagnostik yang bersifat klasifikasi diskrit, estimasi RUL
adalah masalah regresi kontinu yang membutuhkan pemodelan dinamika degradasi
lintas seluruh siklus hidup komponen.

\textcolor{blue}{\subsection{Konstruksi \emph{Health Indicator} dan Normalisasi Degradasi}
\label{subsec:bab2_health_indicator}}

\textcolor{blue}{Sebelum model \emph{deep learning} dapat dilatih untuk estimasi RUL, sinyal vibrasi
mentah perlu diubah menjadi representasi degradasi yang terstruktur, yang dikenal
sebagai \emph{health indicator} (HI). HI merangkum kondisi \emph{bearing} menjadi
skalar atau vektor yang berubah secara monoton seiring penurunan kondisi komponen,
sehingga dapat dipetakan ke label RUL secara langsung.}

\textcolor{blue}{Pendekatan berbasis statistik temporal seperti RMS (\emph{root mean square}) dan
kurtosis rolling window merupakan metode yang paling banyak digunakan karena
kesederhanaan dan kemudahan interpretasinya.
Pada dataset PHM2012, kriteria akhir hidup (\emph{end-of-life}, EOL) didefinisikan
sebagai amplitudo RMS vibrasi yang melampaui ambang 20~g; label RUL kemudian
didefinisikan secara linier dari nilai 1 pada awal operasi hingga 0 pada EOL
\citetitb{Nectoux2012PRONOSTIA}.
Pendekatan berbasis autoencoder \citetitb{Malhotra2016LSTMAnomalyDetect}
mengonstruksi HI dari galat rekonstruksi: model dilatih hanya pada data kondisi
normal, sehingga galat rekonstruksi meningkat secara alami seiring degradasi.}

\textcolor{blue}{Pilihan metode normalisasi HI berpengaruh signifikan terhadap kemampuan model
untuk mempelajari laju degradasi yang konsisten lintas \emph{bearing} yang berbeda.
Penelitian ini menggunakan label RUL linier yang dikalibrasi per rekaman
dengan \emph{StandardScaler} yang disesuaikan pada data pelatihan saja,
untuk mencegah kebocoran informasi dari set uji.
Hubungan antara konstruksi HI, skema normalisasi, dan arsitektur \emph{backbone}
RUL yang digunakan dalam penelitian ini diuraikan lebih lanjut pada §III.4.}

\subsection{Pendekatan Berbasis Fisika dan \emph{Data-Driven}}
\label{subsec:bab2_fisika_datadriven}

Pendekatan berbasis fisika memodelkan mekanisme kerusakan melalui persamaan
diferensial yang parameternya diidentifikasi dari data kondisi. Hukum Paris
untuk perambatan retak dan model keausan Archard memberikan landasan teoretis
yang kuat, namun memerlukan pengetahuan mekanisme fisika yang akurat dan
spesifik per jenis komponen; model yang dikalibrasi untuk satu jenis \emph{bearing}
tidak dapat langsung diterapkan ke \emph{bearing} dengan geometri berbeda
\citetitb{Saxena2008Damage}. Varian hibrida seperti \emph{particle filter}
menggabungkan model fisika dengan pembaruan statistik \emph{online}
\citetitb{An2013ParticleFilter}, namun tetap membutuhkan model struktural yang
benar.

Pendekatan \emph{data-driven} belajar pola degradasi langsung dari data
historis tanpa mengasumsikan model fisika eksplisit. ARIMA
\citetitb{Box1976ARIMA} dan model \emph{proportional hazards}
\citetitb{Cox1972Hazards} meletakkan dasar statistik, namun keduanya
mengasumsikan stasioneritas dan proporsionalitas yang sering tidak terpenuhi
pada sinyal degradasi nonlinier. \emph{Support Vector Regression}
\citetitb{Smola2004SVR} dan \emph{Relevance Vector Machine}
\citetitb{Tipping2001RVM} memberikan alternatif non-parametrik yang populer
sebelum era \emph{deep learning}.

Arsitektur \emph{deep learning} menggeser paradigma ke pembelajaran representasi
hirarkis. LSTM \citetitb{Hochreiter1997LSTM} merupakan arsitektur rekuren
dominan dalam prognostik \emph{bearing} karena mekanisme \emph{gating}-nya untuk
mempertahankan dependensi jangka panjang; BiLSTM dengan mekanisme
\emph{attention} \citetitb{Chen2022BiLSTMAttention} dan TCN dengan konvolusi
kausal berlapis \citetitb{Bai2018TCN} memperluas kemampuan ini. Arsitektur
Transformer \citetitb{Vaswani2017Attention} memberikan fleksibilitas tangkap
konteks global, namun kompleksitasnya kuadratik $\mathcal{O}(L^2)$ terhadap
panjang sekuens menjadi hambatan pada sinyal degradasi yang panjang
\citetitb{Tay2022LRA}.

Perkembangan terbaru berpusat pada dua lini paralel yang mengatasi keterbatasan
Transformer dari sudut berbeda. \emph{State Space Models} (SSM), diwakili
oleh Mamba \citetitb{Gu2023Mamba} dengan parameter selektif bergantung masukan
dan kompleksitas $\mathcal{O}(L)$, memberikan efisiensi komputasi superior
untuk sekuens panjang. xLSTM \citetitb{Beck2024xLSTM} merevitalisasi LSTM
dengan \emph{exponential gating} yang menggantikan saturasi sigmoid dan memori
matriks yang memperluas kapasitas representasi.
\textcolor{blue}{Di lini interpretabilitas berbasis arsitektur, kerangka N-BEATS
\citetitb{Oreshkin2020NBEATS} memperkenalkan dekomposisi basis-blok eksplisit
(\emph{trend} dan \emph{seasonality}) yang menghasilkan struktur internal yang
dapat diinterpretasi tanpa metode XAI pasca-hoc.}
Penelitian ini mengombinasikan
\textcolor{blue}{ketiga keunggulan tersebut --- efisiensi Mamba, kapasitas memori
xLSTM, dan dekomposisi interpretabel N-BEATS ---} dalam tiga arsitektur
\emph{backbone} RUL (Mamba-xLSTM-Net, N-BEATS-xLSTM-RUL, dan SparseGate-TCN-RUL;
diuraikan pada §V.1).

\textbf{Tren terkini.} Peningkatan akurasi prediksi RUL pada dataset
\emph{benchmark} mulai mendekati saturasi: beberapa arsitektur terkini sudah
mencapai RMSE di bawah 0,25 pada PHM2012. Fokus penelitian karenanya bergeser
ke pertanyaan \emph{trustworthiness}: apakah model yang akurat juga
merepresentasikan fisika degradasi secara internal? Pertanyaan inilah yang
dijawab oleh Bab~V melalui pemetaan SAE → BPFx
\citetitb{TotoSuharto2025Journal2RUL}.

\textcolor{blue}{\textbf{Kuantifikasi ketidakpastian (\emph{uncertainty quantification}, UQ).}
Selain akurasi titik estimasi, tren mutakhir dalam prognostik RUL turut menekankan
pentingnya kuantifikasi ketidakpastian: sejauh mana model dapat memberikan
interval kepercayaan yang terkalibrasi, bukan sekadar nilai prediksi tunggal.
Metode seperti \emph{Monte Carlo Dropout} \citetitb{Gal2016MCDropout},
\emph{deep ensemble} \citetitb{Lakshminarayanan2017Ensembles}, dan
\emph{conformal prediction} \citetitb{Angelopoulos2022ConformalSurvey}
mulai banyak diadopsi dalam prognostik mesin rotasi.
Penelitian ini tidak mencakup UQ untuk prediksi RUL secara langsung;
kuantifikasi ketidakpastian yang digunakan — \emph{bootstrap} CI 95\% dan
\emph{permutation test} — diterapkan pada prosedur pemetaan SAE$\to$BPFx
(Bab~V), bukan pada distribusi prediksi RUL itu sendiri.
Perluasan ke UQ prediksi RUL ditetapkan sebagai \emph{future work} pada §VI.5.}

\subsection{Dataset \emph{Benchmark} RUL: PHM2012, XJTU-SY, dan IMS}
\label{subsec:bab2_dataset_rul}

Tiga dataset publik digunakan sebagai \emph{benchmark} dalam penelitian ini.
Ketiganya memiliki karakteristik berbeda, mencakup jenis \emph{bearing}, kondisi
operasi, dan mekanisme kegagalan dominan, sehingga bersama-sama dapat menguji
universalitas pendekatan yang diusulkan.

\textbf{PHM2012 (FEMTO-PRONOSTIA).} Dataset hasil tantangan PHMAP~2012 yang
dikurasi oleh \citenameitb{Nectoux2012PRONOSTIA}. Berisi 17 \emph{bearing} yang
diuji hingga kegagalan di bawah tiga kondisi beban (Kondisi~1: 1.800~rpm,
4.000~N; Kondisi~2: 1.650~rpm, 4.200~N; Kondisi~3: 1.500~rpm, 5.000~N).
Kriteria akhir hidup (EOL): amplitudo RMS vibrasi melampaui ambang 20~g.
Label RUL linier dari 1 (awal operasi) ke 0 (EOL). Dataset dibagi 6 \emph{bearing}
latih dan 11 \emph{bearing} uji (dengan RUL yang dipotong pada waktu tertentu).

\textbf{XJTU-SY.} Dataset dari Xi'an Jiaotong University yang dikurasi oleh
\citenameitb{Wang2020XJTU}. Berisi 10 \emph{bearing} LDK~UER204 yang diuji di bawah
lima kondisi operasi (kombinasi kecepatan dan beban). Jenis kegagalan dominan
adalah \emph{outer-race spalling}, menjadikan dataset ini ideal untuk menguji
apakah representasi laten model menangkap sinyal BPFO.

\textbf{IMS (Universitas Cincinnati).} Dataset yang diperkenalkan oleh
\citenameitb{Qiu2006WaveletBearing}, berisi empat \emph{bearing} Rexnord yang
dijalankan secara bersamaan pada kecepatan 2.000~rpm dengan akuisisi tiap
10~menit. \emph{Bearing}~3 mengalami kegagalan \emph{outer-race}; \emph{Bearing}~4
mengalami kelelahan elemen bergulir (\emph{rolling-element fatigue}).
Dengan hanya empat \emph{bearing}, dataset ini lebih menantang secara statistik
dan menguji kemampuan generalisasi pada data berskala kecil.

Tabel~\ref{tab:bab2_dataset_rul} merangkum karakteristik ketiga dataset.
Dataset CWRU juga digunakan secara tambahan pada Bab~V sebagai validasi
lintas dataset untuk SAE-BPFx, namun dengan catatan penting bahwa $n=10$
rekaman menjadikannya \emph{underpowered} secara statistik untuk pengujian
korelasi Pearson (p-value > 0,05 meski nilai \emph{hit-rate} tinggi).

\begin{table}[ht]
  \centering
  \caption{Karakteristik tiga dataset \emph{benchmark} prediksi sisa umur
           pakai \emph{rolling element bearing} yang digunakan dalam penelitian ini.}
  \label{tab:bab2_dataset_rul}
  \footnotesize
  \begin{tabularx}{\textwidth}{@{}l l l c l c@{}}
    \toprule
    \textbf{Dataset} & \textbf{Institusi} & \textbf{Tipe \emph{bearing}}
      & \textbf{Jumlah} & \textbf{Kondisi operasi} & \textbf{Kriteria EOL} \\
    \midrule
    PHM2012 & FEMTO-PRONOSTIA & SKF/NTN & 17 & 3 kondisi beban & RMS $>$ 20~g \\
    XJTU-SY & Xi'an Jiaotong Univ. & LDK UER204 & 10 & 5 kondisi & Sinyal berhenti \\
    IMS     & Univ. Cincinnati & Rexnord & 4  & 1 kondisi (2.000~rpm) & Label manual \\
    \bottomrule
  \end{tabularx}

  \smallskip
  {\footnotesize Sumber: \citenameitb{Nectoux2012PRONOSTIA} (PHM2012);
   \citenameitb{Wang2020XJTU} (XJTU-SY); \citenameitb{Qiu2006WaveletBearing}
   (IMS).}
\end{table}

% -----------------------------------------------------------------
\section{Interpretabilitas Mekanistik dan \emph{Sparse Autoencoder}}
\label{sec:bab2_mekanistik}
% -----------------------------------------------------------------

Metode XAI yang dibahas pada §\ref{sec:bab2_xai} beroperasi di tingkat
\emph{masukan}: mengidentifikasi fitur atau posisi sinyal yang paling berpengaruh
terhadap prediksi. Pertanyaan yang lebih dalam, yaitu \emph{konsep fisika apa yang
direpresentasikan model pada lapisan tersembunyinya?}, tidak dapat dijawab
oleh SHAP atau metode atribusi masukan lainnya karena keduanya tidak beroperasi
pada ruang representasi laten. Interpretabilitas mekanistik mengisi celah ini
dengan membedah representasi internal model.

\subsection{Hipotesis Superposisi dan Top-\textit{k} SAE}
\label{subsec:bab2_superposisi_sae}

\citenameitb{Elhage2022Superposition} memformulasikan \emph{hipotesis
superposisi}: jaringan saraf merepresentasikan lebih banyak fitur daripada
jumlah dimensi latennya dengan cara menyimpan fitur-fitur dalam pola superposisi
yang tidak berinterferensi. Akibatnya, setiap neuron cenderung merespons
terhadap banyak konsep yang tidak berkaitan (\emph{polysemantik}); perilaku ini
menyulitkan interpretabilitas mekanistik.

\emph{Sparse Autoencoder} (SAE) \citetitb{Makhzani2014kSAE} menawarkan solusi:
melatih autoencoder dengan kendala \emph{sparsity} pada representasi laten
sehingga setiap fitur SAE yang aktif cenderung merespons satu konsep tunggal
(\emph{monosemantik}). \citenameitb{Bricken2023Monosemanticity} mendemonstrasikan
secara empiris pada model bahasa kecil bahwa SAE mampu mengekstrak fitur-fitur
yang berkorespondensi dengan konsep yang dapat diinterpretasi seperti nama
orang, bahasa pemrograman, dan emosi.

Varian Top-$k$ SAE \citetitb{Cunningham2023TopkSAE} menggunakan strategi aktivasi
eksplisit: pada setiap vektor laten, hanya $k$ fitur dengan nilai terbesar yang
dipertahankan aktif, sisanya di-nol-kan. Dibandingkan dengan regularisasi L1
konvensional, Top-$k$ memberikan kendali langsung atas tingkat \emph{sparsity}
dan pelatihan yang lebih stabil. \citenameitb{Templeton2024Scaling} mengvalidasi
pendekatan ini pada model produksi skala besar, menunjukkan bahwa jutaan fitur
yang dapat diinterpretasi dapat diekstrak.

\subsection{Gap: Belum Ada SAE untuk Domain Prognostik \emph{Bearing}}
\label{subsec:bab2_gap_sae}

Seluruh penelitian interpretabilitas mekanistik berbasis SAE hingga saat ini
berkonsentrasi pada model bahasa besar. Domain prognostik \emph{bearing} menawarkan
kondisi yang justru lebih ideal sebagai \emph{test case}: konsep fisika yang
harus direpresentasikan oleh model RUL yang baik, yaitu frekuensi karakteristik
BPFO, BPFI, BSF, dan FTF, terdefinisi secara matematis dan terukur secara
langsung dari sinyal vibrasi. Kondisi ini sangat berbeda dari domain bahasa, yang interpretasi
\emph{ground-truth}-nya bersifat subjektif dan sulit diverifikasi secara
kuantitatif.

Terdapat hipotesis fisika yang kuat: model RUL yang dilatih pada data degradasi
nyata \emph{seharusnya} menginternalisasi informasi BPFx dalam ruang latennya,
karena frekuensi-frekuensi ini merupakan tanda fisik kunci dari mekanisme
kegagalan yang dipelajari. Namun, belum ada studi yang mencoba memetakan fitur
SAE dari representasi laten model RUL ke BPFx secara statistik signifikan,
dengan kontrol terhadap artefak arsitektur dan prosedur.

Penelitian yang dilaporkan pada \citetitb{TotoSuharto2025Journal2RUL}
merupakan studi pertama yang: (i)~mengadaptasi Top-$k$ SAE ke domain
prognostik \emph{bearing} (\emph{novelti N3: SAEBearing}), (ii)~memetakan
fitur SAE ke BPFx menggunakan korelasi Pearson dan uji statistik \emph{bootstrap}
CI~95\% \textcolor{blue}{($B = 1.000$ resampel)} + \emph{permutation test}
\textcolor{blue}{dua-sisi} dengan dua \emph{negative controls} (model
Xavier-init dan \emph{hidden state} Gaussian acak), dan (iii)~mengkonfirmasi
bahwa korespondensi SAE $\leftrightarrow$ BPFx merupakan properti \emph{emergent}
dari representasi yang terlatih, bukan artefak arsitektur atau prosedur.

% -----------------------------------------------------------------
% §II.7 BARU — tinjauan XAI lintas domain sinyal 1-D
% -----------------------------------------------------------------
\textcolor{blue}{\section{XAI pada Domain Sinyal Satu Dimensi: Tinjauan Lintas Domain}
\label{sec:bab2_xai_lintas_domain}}

\textcolor{blue}{Penerapan XAI pada sinyal satu dimensi tidak terbatas pada domain vibrasi mesin.
Tinjauan lintas domain memberikan konteks penting untuk menilai sejauh mana FSM
merupakan kontribusi metodologis yang bersifat umum, dan di mana celah spesifik
domain vibrasi yang belum terisi.}

\textcolor{blue}{Pada domain \emph{electrocardiogram} (ECG), SHAP dan \emph{Integrated Gradients}
(IG) telah digunakan untuk mengidentifikasi segmen temporal dalam denyut jantung
yang paling berpengaruh terhadap klasifikasi aritmia.
\citenameitb{Strodthoff2021ECGExplainability} menerapkan beberapa metode XAI
secara komparatif pada dataset ECG publik dan melaporkan bahwa SHAP
\emph{DeepExplainer} menghasilkan atribusi yang lebih stabil dibanding
\emph{saliency maps} vanila, meskipun sensitivitas terhadap pilihan
\emph{background} tetap menjadi kendala.
Pada domain audio dan \emph{speech}, Grad-CAM yang diadaptasi untuk konvolusi
1-D \citetitb{Ribeiro2016LIME} dan LIME berbasis segmen temporal telah digunakan
untuk mengidentifikasi bagian ucapan yang menentukan klasifikasi emosi atau
kata kunci.}

\textcolor{blue}{Kesamaan lintas domain ini menegaskan bahwa pendekatan XAI berbasis atribusi
posisi temporal---yang menjadi fondasi FSM---bukan merupakan pendekatan yang
asing. Namun, perbedaan mendasar antara domain ECG/audio dan domain vibrasi
\emph{bearing} terletak pada interpretabilitas fisik: sinyal ECG dan audio
tidak memiliki frekuensi karakteristik terdefinisi matematis yang setara
dengan BPFO/BPFI/BSF/FTF, sehingga validasi fisika atribusi (sebagaimana
dilakukan pada FSM dan SAE-BPFx) tidak dapat dilakukan secara kuantitatif
di domain tersebut.
Di domain vibrasi \emph{bearing} sendiri, seluruh studi XAI yang ada beroperasi
di tingkat fitur agregat (RMS, kurtosis, energi pita spektral), bukan pada
resolusi sinyal mentah penuh.
FSM mengisi celah ini sebagai kerangka atribusi sinyal mentah resolusi penuh
pertama untuk WDCNN pada dataset CWRU \citetitb{TotoSuharto2025Journal1FSM}.}

% -----------------------------------------------------------------
\section{Peta Kesenjangan Literatur dan Posisi Penelitian}
\label{sec:bab2_gap_literatur}
% -----------------------------------------------------------------

\textcolor{blue}{Berdasarkan tinjauan pada §\ref{sec:bab2_pdm_industri}
hingga §\ref{sec:bab2_xai_lintas_domain}, bagian ini memetakan
\emph{state of the art} penelitian perawatan prediktif secara komprehensif
dan mengidentifikasi kesenjangan yang menjadi justifikasi langsung bagi
ketiga rumusan masalah pada §\ref{sec:bab1_rumusan}.}

\textcolor{blue}{Tabel~\ref{tab:bab2_sota_pdm} menyajikan pemaparan detail
\emph{state of the art} perawatan prediktif \emph{bearing} dari 17 studi
representatif yang dikelompokkan berdasarkan lima tema utama, plus empat
karya internal disertasi ini. Setiap baris memuat metode, dataset, hasil
utama, dan keterbatasan masing-masing studi.}

%% =====================================================================
%% TABEL III-1: SotA PdM (longtable, 5 grup tematik + posisi disertasi)
%% =====================================================================
\begin{small}
\setlength{\tabcolsep}{3pt}
\renewcommand{\arraystretch}{1.15}
\begin{longtable}{@{}>{\raggedright\arraybackslash}p{2.0cm}
                     >{\raggedright\arraybackslash}p{3.0cm}
                     >{\raggedright\arraybackslash}p{2.4cm}
                     >{\raggedright\arraybackslash}p{7.8cm}@{}}
\caption{\textcolor{blue}{Pemaparan \emph{state of the art} (SotA) penelitian
         perawatan prediktif (PdM) \emph{bearing} dan posisi penelitian
         disertasi ini, dikelompokkan dalam lima grup tematik (A--E) ditambah
         karya internal (Grup F).}}
\label{tab:bab2_sota_pdm} \\
\toprule
\textbf{Studi} & \textbf{Metode} & \textbf{Dataset} &
\textbf{Hasil Utama dan Keterbatasan} \\
\midrule
\endfirsthead

\multicolumn{4}{l}{\footnotesize\itshape \tablename~\thetable\ (lanjutan)} \\
\toprule
\textbf{Studi} & \textbf{Metode} & \textbf{Dataset} &
\textbf{Hasil Utama dan Keterbatasan} \\
\midrule
\endhead

\midrule
\multicolumn{4}{r}{\footnotesize\itshape (bersambung pada halaman berikutnya)} \\
\endfoot

\bottomrule
\endlastfoot

\multicolumn{4}{@{}l}{\textbf{\textcolor{blue}{Grup A. Diagnostik
\emph{bearing} berbasis pembelajaran mesin dan pembelajaran mendalam}}} \\
\midrule

\citenameitb{Smith2015CWRU} &
Diagnosa \emph{bearing} CWRU dengan ekstraksi fitur klasik (RMS, kurtosis,
\emph{envelope spectrum}) &
CWRU (12~kHz/48~kHz, 10 kelas) &
Memvalidasi CWRU sebagai \emph{benchmark} standar diagnostik
\emph{bearing}; namun bergantung pada fitur rekayasa manual sehingga
sulit digeneralisasi ke kondisi operasi baru. \\[3pt]

\citenameitb{Zhang2017WDCNN} &
\emph{Wide Deep CNN} (WDCNN) dengan kernel konvolusi lebar
(lebar 64, \emph{stride} 16) langsung pada sinyal mentah &
CWRU (segmen 2.048 titik) &
Akurasi $>$~99\,\% pada CWRU tanpa ekstraksi fitur manual.
Keterbatasan: tidak menyertakan analisis interpretabilitas;
\emph{black-box} bagi insinyur perawatan. \\[3pt]

\citenameitb{Ayvaz2021PdMML} &
Gabungan ANN + SVM untuk klasifikasi kondisi mesin &
Data sensor produksi otomotif &
Mendemonstrasikan kelayakan ML pada lingkungan produksi; namun
hanya membandingkan dua algoritma dan tidak menggunakan protokol
evaluasi temporal. \\[3pt]

\citenameitb{Cui2022AnomalyML} &
SVM + k-NN untuk deteksi anomali multi-mode &
\emph{Bearing} eksperimental (\emph{rig} laboratorium) &
\emph{Recall} tinggi pada deteksi anomali jarang; tetapi
tidak ada studi ablasi atau \emph{cross-dataset validation}. \\[3pt]

\citenameitb{Marani2020ANFIS} &
ANFIS (\emph{Adaptive Neuro-Fuzzy Inference System}) berbasis fitur
statistik &
Dataset \emph{bearing} laboratorium &
Interpretabilitas linguistik melalui aturan fuzzy; namun
\emph{scaling} ke ribuan kelas dan sinyal panjang masih bermasalah
dan akurasi tertinggal dibandingkan DL modern. \\

\midrule
\multicolumn{4}{@{}l}{\textbf{\textcolor{blue}{Grup B. \emph{Explainable AI}
(XAI) pada diagnostik \emph{bearing}}}} \\
\midrule

\citenameitb{Santos2024SHAPSVM} &
SHAP \emph{KernelExplainer} pada 9--15 fitur statistik dengan model SVM &
\emph{Bearing} CWRU (subset) &
Memberikan peringkat kepentingan fitur secara global; tetapi
beroperasi pada level fitur agregat (kurtosis, \emph{crest factor},
energi pita), bukan pada sinyal mentah. \\[3pt]

\citenameitb{Jeon2023CAMBearing} &
\emph{Class Activation Mapping} (CAM) pada spektrogram STFT &
\emph{Bearing} laboratorium &
Visualisasi area frekuensi--waktu yang menentukan klasifikasi;
namun beroperasi pada representasi spektrogram, bukan pada
sinyal mentah resolusi penuh; resolusi spektrogram membatasi
lokalisasi temporal halus. \\[3pt]

\citenameitb{Bengherbia2024SHAPBearing} &
SHAP untuk diagnostik \emph{bearing} pada model XGBoost dan ANN &
\emph{Bearing} CWRU + dataset internal &
Validasi tingkat fitur agregat; tidak meninjau apakah model
benar-benar belajar mekanisme fisik impuls transien BPFx. \\[3pt]

\citenameitb{Serradilla2022DLPdM} &
Survei XAI untuk PdM berbasis DL: SHAP, IG, \emph{attention rollout} &
Meta-analisis 60+ studi &
Memetakan lanskap XAI untuk PdM; menemukan bahwa hampir semua
studi beroperasi pada tingkat masukan dan tidak pada
representasi laten model. \\

\midrule
\multicolumn{4}{@{}l}{\textbf{\textcolor{blue}{Grup C. Prognostik
\emph{bearing} dan estimasi sisa umur pakai (RUL)}}} \\
\midrule

\citenameitb{Hochreiter1997LSTM} &
\emph{Long Short-Term Memory} (LSTM) untuk pemodelan sekuens
temporal panjang &
Tugas sekuens sintetis &
Fondasi prediksi temporal jangka panjang; namun rentan
\emph{vanishing gradient} pada sekuens sangat panjang
($>$~10.000 titik) yang lazim pada degradasi \emph{bearing}. \\[3pt]

\citenameitb{Vaswani2017Attention} &
\emph{Transformer} dengan \emph{self-attention} murni &
WMT (NLP) &
Pemodelan dependensi jangka panjang tanpa rekurens; namun
kompleksitas $O(n^2)$ menjadi mahal untuk sinyal vibrasi panjang. \\[3pt]

\citenameitb{Gu2023Mamba} &
\emph{State Space Model} (SSM) selektif dengan parameter bergantung
masukan &
Long-Range Arena, NLP, audio &
Kompleksitas linear $O(n)$ dengan kinerja menyamai Transformer;
diadaptasi dalam disertasi ini sebagai \emph{backbone} prognostik. \\[3pt]

\citenameitb{Beck2024xLSTM} &
\emph{Extended LSTM} dengan \emph{exponential gating} dan
\emph{matrix memory} &
Bahasa, audio, sekuens panjang &
Mengatasi keterbatasan LSTM klasik; namun belum pernah diterapkan
pada prognostik \emph{bearing} sebelum disertasi ini. \\[3pt]

\citenameitb{Khorram2020CNNLSTM} &
CNN-LSTM \emph{hybrid} untuk estimasi RUL &
PHM2012 (FEMTO-PRONOSTIA) &
RMSE kompetitif pada PHM2012; namun tanpa XAI dan tanpa validasi
lintas \emph{rig} eksperimental. \\[3pt]

\citenameitb{Lei2018PrognosticsReview} &
Tinjauan komprehensif metode prognostik berbasis data dan berbasis
fisika &
Meta-analisis &
Memetakan dua paradigma; menyoroti bahwa model \emph{data-driven}
yang akurat sering kali tidak dapat diinterpretasikan secara fisika. \\

\midrule
\multicolumn{4}{@{}l}{\textbf{\textcolor{blue}{Grup D. Interpretabilitas
mekanistik dengan \emph{Sparse Autoencoder} (SAE)}}} \\
\midrule

\citenameitb{Bricken2023Monosemanticity} &
SAE pada lapisan MLP \emph{transformer} untuk ekstraksi fitur
monosemantik &
Model bahasa kecil (Anthropic) &
Membuktikan bahwa SAE mengekstrak konsep yang dapat dipetakan
ke konsep semantik; namun domain bahasa, belum domain sinyal fisik. \\[3pt]

\citenameitb{Cunningham2023TopkSAE} &
\emph{Top-k Sparse Autoencoder} dengan \emph{sparsity constraint}
yang dapat dikendalikan &
Model bahasa &
Memberikan kontrol langsung atas sparsitas; diadaptasi dalam
disertasi ini ke domain prognostik \emph{bearing}. \\[3pt]

\citenameitb{Templeton2024Scaling} &
\emph{Scaling} SAE pada \emph{Claude 3 Sonnet} hingga jutaan fitur &
LLM produksi &
Membuktikan SAE dapat diperluas; namun belum diadaptasi ke domain
non-bahasa dengan validasi fisika sebelum disertasi ini. \\

\midrule
\multicolumn{4}{@{}l}{\textbf{\textcolor{blue}{Grup E. Arsitektur
sistem PdM \emph{edge}/\emph{cloud}}}} \\
\midrule

\citenameitb{Cheng2019EdgeCloud} &
Koordinasi \emph{edge}--\emph{cloud} untuk pemrosesan data PdM &
Studi kasus industri manufaktur &
Membahas pembagian beban komputasi; namun masih
\emph{cloud-centric} dan tanpa lapisan XAI per \emph{tier}. \\[3pt]

\citenameitb{Hafeez2021EdgePdM} &
\emph{Edge intelligence} untuk PdM dengan model ringan &
Sensor IoT industri &
Mendemonstrasikan kelayakan \emph{edge deployment}; namun
tidak menyertakan lapisan interpretabilitas yang dapat diaudit. \\

\midrule
\multicolumn{4}{@{}l}{\textbf{\textcolor{blue}{Grup F. Karya internal
disertasi ini}}} \\
\midrule

\citenameitb{TotoSuharto2024Conf1SVM} &
SVM-RBF + Regresi Logistik dengan SHAP \emph{KernelExplainer} pada
vektor fitur HI 36-D &
CWRU (\emph{temporal split} 54/13/33\,\%) &
SVM-RBF mencapai akurasi 96,4\,\%; XAI level fitur (Conf~1).
Mendukung Tier Edge IoT pada arsitektur multi-tier. \\[3pt]

\citenameitb{TotoSuharto2024Conf2Tree} &
Decision Tree, Random Forest, XGBoost dengan SHAP \emph{TreeExplainer}
pada vektor fitur HI 36-D &
CWRU (\emph{temporal split} 54/13/33\,\%) &
XGBoost mencapai akurasi 97,1\,\%; XAI level fitur dengan
\emph{tree-path} eksak (Conf~2). Mendukung Tier Edge Server. \\[3pt]

\citenameitb{TotoSuharto2025Journal1FSM} &
WDCNN + SHAP \emph{DeepExplainer} pada sinyal mentah 2.048 titik
$\to$ \emph{Fault Signature Maps} (FSM, 3 varian: Signed, Absolute,
Variance) &
CWRU (1.240/310/750 segmen) &
Akurasi 99,87\,\% (Var-A); \emph{discriminability} $D = 0{,}216$,
\emph{split-half stability} $r = 0{,}940$; ablasi \emph{BatchNorm}
menyingkap \emph{trade-off}: $+133$\,\% diskriminabilitas dengan
$-3{,}60$~pp akurasi (J~1). \\[3pt]

\citenameitb{TotoSuharto2025Journal2RUL} &
Tiga \emph{backbone} RUL (Mamba-xLSTM-Net, N-BEATS-xLSTM-RUL,
SparseGate-TCN-RUL) + \emph{Top-k SAE} dipetakan ke BPFx melalui
korelasi dengan \emph{Hilbert envelope} &
PHM2012, XJTU-SY, IMS, CWRU &
\emph{Hit-rate} BPFI 2,3\,\% (PHM2012, $p < 0{,}001$), BPFO
2,2\,\% (XJTU-SY, $p < 0{,}001$); konsistensi urutan dominansi BPFx
lintas tiga arsitektur; validasi dengan \emph{bootstrap} CI 95\,\%,
\emph{permutation test}, dan dua \emph{negative controls} (J~2). \\

\end{longtable}
\end{small}

\textcolor{blue}{Berdasarkan Tabel~\ref{tab:bab2_sota_pdm}, dapat
diidentifikasi bahwa sebagian besar studi diagnostik \emph{bearing}
(Grup~A) hanya membandingkan satu hingga tiga algoritma tanpa XAI dan tanpa
protokol evaluasi temporal yang ketat. Studi XAI yang ada (Grup~B) beroperasi
pada tingkat fitur agregat atau spektrogram, bukan pada resolusi sinyal
mentah. Di bidang prognostik (Grup~C), arsitektur modern seperti Mamba dan
xLSTM belum diterapkan pada \emph{bearing}, dan interpretabilitas mekanistik
berbasis SAE (Grup~D) baru dikembangkan untuk model bahasa besar dan belum
pernah diadaptasi ke domain mesin rotasi. Arsitektur sistem PdM
\emph{edge}/\emph{cloud} (Grup~E) belum mengintegrasikan XAI per lapis
secara terstruktur. Empat karya internal pada Grup~F mengisi celah-celah
tersebut dengan dukungan empiris dan inferensi statistik.}

Tabel~\ref{tab:bab2_posisi_penelitian} selanjutnya
menyajikan posisi keempat karya penelitian ini relatif terhadap literatur umum pada
dimensi-dimensi yang menutup kesenjangan tersebut.

\begin{table}[ht]
  \centering
  \caption{Posisi empat karya penelitian ini relatif terhadap literatur umum
           pada \textcolor{blue}{enam} dimensi kritis. Simbol \checkmark = terpenuhi,
           \emph{sebagian} = terpenuhi parsial; kolom tidak relevan
           atau tidak dilaporkan dikosongkan.}
  \label{tab:bab2_posisi_penelitian}
  \footnotesize
  \setlength{\tabcolsep}{4pt}
  \renewcommand{\arraystretch}{1.15}
  \begin{tabularx}{\textwidth}{@{}>{\raggedright\arraybackslash}p{3.0cm} *{5}{>{\raggedright\arraybackslash}X}@{}}
    \toprule
    \textbf{Aspek} & \textbf{Conf~1} & \textbf{Conf~2} & \textbf{J~1}
      & \textbf{J~2} & \textbf{Literatur umum} \\
    \midrule
    Level XAI
      & Fitur & Fitur & Sinyal & Laten & Fitur atau tidak ada \\[3pt]
    Dataset diagnostik (CWRU)
      & \checkmark & \checkmark & \checkmark & \emph{sebagian} & Beragam \\[3pt]
    Dataset prognostik (PHM/XJTU/IMS)
      & --- & --- & --- & \checkmark & PHM dan XJTU dominan \\[3pt]
    Validasi statistik
      & Akurasi, F1 & Akurasi, F1 & Akurasi, F1 + 3 metrik FSM
      & RMSE + \emph{bootstrap} CI + uji permutasi & Akurasi saja \\[3pt]
    \textcolor{blue}{Validasi industri (data real)}
      & \textcolor{blue}{$\times$} & \textcolor{blue}{$\times$}
      & \textcolor{blue}{$\times$}
      & \textcolor{blue}{\emph{sebagian} (SKF)}
      & \textcolor{blue}{Sangat terbatas} \\[3pt]
    Integrasi diagnostik–prognostik
      & $\times$ & $\times$ & $\times$ & $\times$ & $\times$ \\[3pt]
    Interpretabilitas dua lapis (masukan + laten)
      & $\times$ & $\times$ & $\times$ (hanya masukan) & $\times$ (hanya laten)
      & $\times$ \\
    \bottomrule
  \end{tabularx}

  \smallskip
  {\footnotesize Conf~1 = \citenameitb{TotoSuharto2024Conf1SVM} (SVM/LR);
   Conf~2 = \citenameitb{TotoSuharto2024Conf2Tree} (DT/RF/XGB);
   J~1 = \citenameitb{TotoSuharto2025Journal1FSM} (WDCNN+FSM);
   J~2 = \citenameitb{TotoSuharto2025Journal2RUL} (RUL+SAE-BPFx).}
\end{table}

Tabel~\ref{tab:bab2_posisi_penelitian} mengungkapkan
\textcolor{blue}{empat} kesenjangan yang
belum terisi oleh literatur yang ada, termasuk oleh keempat karya secara
individual.

\textbf{Kesenjangan~1 (Diagnostik level sinyal).} Hampir seluruh penelitian
XAI pada klasifikasi \emph{bearing} beroperasi di tingkat fitur rekayasa manual:
menjelaskan \emph{fitur mana} (kurtosis, crest factor, dst.) yang berkontribusi,
bukan \emph{posisi mana dalam sinyal mentah 2.048 sampel} yang menjadi basis
keputusan. Belum ada peta fingerprint kelas kerusakan berbasis atribusi sinyal
mentah dengan validasi multi-metrik (discriminability, severity monotonicity,
split-half stability). Kesenjangan ini dijawab oleh Journal~1 dan Bab~IV
\textcolor{blue}{(RM-2)}.

\textbf{Kesenjangan~2 (Interpretabilitas laten prognostik).} Belum ada studi
yang memetakan fitur SAE dari representasi laten model RUL ke frekuensi
karakteristik fisik \emph{bearing} secara statistik signifikan, dengan kontrol
terhadap artefak arsitektur (model Xavier-init) dan artefak prosedur (hidden
state acak). Metode interpretabilitas yang ada beroperasi di tingkat masukan atau
tidak dapat menjembatani representasi laten dengan teori fisika degradasi.
Kesenjangan ini dijawab oleh Journal~2 dan Bab~V
\textcolor{blue}{(RM-2)}.

\textbf{Kesenjangan~3 (Integrasi diagnostik–prognostik yang dapat diaudit).}
Belum ada kerangka yang menyatukan interpretabilitas tingkat masukan (FSM untuk
diagnostik) dan interpretabilitas tingkat laten (SAE-BPFx untuk prognostik)
menjadi satu sistem PdM multi-tier yang dapat diaudit lapis demi lapis.
Ketiadaan integrasi ini merupakan hambatan adopsi industri: teknisi tidak dapat
menelusuri alasan keputusan dari sinyal mentah hingga konsep fisika laten.
Kesenjangan ini dijawab oleh \textcolor{blue}{arsitektur multi-tier pada Bab~VI
(RM-3)}.

\textcolor{blue}{\textbf{Kesenjangan~4 (Validasi industri real).} Sebagian
besar studi diagnostik dan prognostik \emph{bearing} dikembangkan dan
divalidasi pada dataset publik yang dikumpulkan dalam kondisi laboratorium
terkontrol. Validasi pada lingkungan industri nyata, khususnya pada proses
produksi \emph{ball bearing} dengan kerusakan yang berkembang secara alami,
masih sangat terbatas. Model yang akurat pada dataset benchmark belum tentu
bekerja pada data industri real. Kesenjangan ini dijawab oleh validasi RUL
pada data vibrasi PT~SKF Indonesia (RM-3).}

\textcolor{blue}{Keempat} kesenjangan tersebut terhubung langsung dengan
\textcolor{blue}{tiga} pertanyaan penelitian
pada §\ref{sec:bab1_rumusan}: \textcolor{blue}{Kesenjangan~1 dan~2
menjustifikasi RM-2 (XAI dua lapis); Kesenjangan~3 dan~4 menjustifikasi
RM-3 (validasi industri dan arsitektur multi-tier). RM-1 (model PdM
diagnostik dan prognostik yang akurat) merupakan prasyarat yang memungkinkan
seluruh kesenjangan di atas ditangani.} Metodologi umum untuk menjawab
\textcolor{blue}{keempat} kesenjangan disajikan pada Bab~III, dengan detail
diagnostik pada Bab~IV dan prognostik pada Bab~V.

% =====================================================================
% CATATAN REFERENSI BARU (perlu ditambahkan ke daftar pustaka)
% Tandai dengan [BARU] untuk kemudahan pelacakan
% =====================================================================
%
% [BARU-1] Li, X., Zhang, W., Ding, Q., & Sun, J.-Q. (2020).
%   Intelligent rotating machinery fault diagnosis based on deep learning
%   using data augmentation. Journal of Intelligent Manufacturing, 31(2),
%   433--452.
%   -> Digunakan untuk: \citetitb{Li2020TransferReview}
%   -> Verifikasi: https://doi.org/10.1007/s10845-018-1456-1
%
% [BARU-2] Ganin, Y., & Lempitsky, V. (2015).
%   Unsupervised domain adaptation by backpropagation.
%   ICML 2015, PMLR 37, 1180--1189.
%   -> Digunakan untuk: \citetitb{Ganin2015DANN}
%   -> Verifikasi: https://proceedings.mlr.press/v37/ganin15.html
%
% [BARU-3] Malhotra, P., Ramakrishnan, A., Anand, G., Vig, L., Agarwal, P.,
%   & Shroff, G. (2016). LSTM-based encoder-decoder for multi-sensor
%   anomaly detection. arXiv:1607.00148.
%   -> Digunakan untuk: \citetitb{Malhotra2016LSTMAnomalyDetect}
%   -> Verifikasi: https://arxiv.org/abs/1607.00148
%
% [BARU-4] Gal, Y., & Ghahramani, Z. (2016). Dropout as a Bayesian
%   approximation: Representing model uncertainty in deep learning.
%   ICML 2016, PMLR 48, 1050--1059.
%   -> Digunakan untuk: \citetitb{Gal2016MCDropout}
%   -> Verifikasi: https://proceedings.mlr.press/v48/gal16.html
%
% [BARU-5] Lakshminarayanan, B., Pritzel, A., & Blundell, C. (2017).
%   Simple and scalable predictive uncertainty estimation using deep
%   ensembles. NeurIPS 2017, 6402--6413.
%   -> Digunakan untuk: \citetitb{Lakshminarayanan2017Ensembles}
%   -> Verifikasi: https://papers.nips.cc/paper/2017/hash/...
%
% [BARU-6] Angelopoulos, A. N., & Bates, S. (2022). A gentle introduction
%   to conformal prediction and distribution-free uncertainty quantification.
%   arXiv:2107.07511.
%   -> Digunakan untuk: \citetitb{Angelopoulos2022ConformalSurvey}
%   -> Verifikasi: https://arxiv.org/abs/2107.07511
%
% [BARU-7] Strodthoff, N., Wagner, P., Schaeffter, T., & Samek, W. (2021).
%   Deep learning for ECG analysis: Benchmarks and insights from PTB-XL.
%   IEEE Journal of Biomedical and Health Informatics, 25(5), 1519--1528.
%   -> Digunakan untuk: \citetitb{Strodthoff2021ECGExplainability}
%   -> Verifikasi: https://doi.org/10.1109/JBHI.2020.3022989
%
% Referensi baris tabel (perlu verifikasi manual):
% [BARU-T1] Santos dkk. (2021) -- SVM CWRU 95.1% + KernelExplainer
%   Placeholder; cek: Santos, P., et al. (2021). IEEE Access atau sejenis.
% [BARU-T2] Cerrada dkk. (2022) -- XGBoost CWRU 96.8%
%   Placeholder; cek: Cerrada, M., et al. (2022). Eksak judul perlu verifikasi.
% =====================================================================

% =====================================================================
% AUDIT KONSISTENSI REFERENSI & ANGKA (lintas bab) — putaran ini
% =====================================================================
% A. REFERENSI
%   - 51 referensi asli Bab II: SEMUA dipertahankan (tidak ada yang hilang).
%   - 7 referensi baru (biru): Li2020TransferReview, Ganin2015DANN,
%     Malhotra2016LSTMAnomalyDetect, Gal2016MCDropout,
%     Lakshminarayanan2017Ensembles, Angelopoulos2022ConformalSurvey,
%     Strodthoff2021ECGExplainability.
%   - DITAMBAHKAN putaran ini: Oreshkin2020NBEATS (dipakai Bab V tapi belum
%     ada di Bab II) -> kini disitir di §II.5.2 (biru).
%   - Loshchilov2019AdamW (dipakai Bab V) sengaja TIDAK dimasukkan ke Bab II:
%     bersifat detail implementasi optimizer, bukan landasan teori.
%
% B. KOREKSI ANGKA agar konsisten dgn sumber otoritatif (Bab IV / Journal FSM):
%   1. BPFx CWRU: 9 bola/0deg/29,95Hz/107,36;162,18;70,87;11,93
%      -> 8 bola/15,17deg/1.750rpm/104,6;157,9;68,9;13,1  (sesuai §IV.11)
%   2. WDCNN Var C akurasi: 96,27% -> 96,13% (sesuai Tabel IV CV)
%   3. Stabilitas FSM Var A: ANGKA OTORITATIF = 0,940 (analisis utama,
%      300 bg/500 uji; SAMA dengan kolom Stabilitas Tabel IV & Journal FSM
%      "split-half stability = 0.940, 94% cosine"). Catatan lama yg menyebut
%      0,940 KELIRU/seharusnya 0,873 -> DIBATALKAN: 0,873 tidak ada di sumber.
%      Pada subset ablasi (30 bg/50 uji): Var A stab 0,705 -> Var C stab 0,741.
%   4. Klaim "+133%": pasangan yg BENAR adalah 0,315 -> 0,735 (subset ablasi),
%      = +0,420 poin = 2,3x. BUKAN 0,216->0,735 (itu mencampur metric-set
%      utama dgn ablasi). 0,216 adalah diskriminabilitas analisis utama Var A.
%   5. Permutation test: definisi "proporsi |r_null| > |r_obs|" adalah
%      DUA-SISI (sesuai Journal RUL: "two-sided p-value ... absolute
%      correlation"). Bab II sudah benar memakai "dua-sisi"; Bab V baris 337/342
%      yg memakai "satu-sisi (greater)" KELIRU dan telah dikoreksi ke "dua-sisi".
%
% C. CATATAN LINTAS-DOKUMEN (di luar Bab II, perlu tindak lanjut terpisah):
%   - Permutasi B=1.000: SUDAH selaras di Pendahuluan (§I, B=1.000), Bab V,
%     Bab VI, dan Bab II. Konflik lama n=10.000 sudah tidak ada.
%   - Arah uji permutasi: Bab V DIKOREKSI dari "satu-sisi" -> "dua-sisi"
%     agar konsisten dgn definisi |r_null|>|r_obs| dan Journal RUL.
%   - Bab VI DIKOREKSI: pasangan +133% diselaraskan ke 0,315->0,735 (ablasi),
%     dan framing "BN menekan 133%" diubah ke "penghapusan BN menaikkan 133%".
%   - Stabilitas FSM: Bab I & Bab II DIKOREKSI ke 0,940 (analisis utama),
%     menggantikan 0,873/0,891 yg tidak bersumber.
% =====================================================================
